\subsection{直线与闭超平面}

命题 2. --- 非离散赋值除环 K 上的每个一维 Hausdorff 拓扑向量空间 E 都同构于 \(K_{s}\)；事实上，对 E 中每个 \(a \neq 0\)，映射 \(\xi \mapsto \xi a\)（\(K_{s}\) 到 E 上的映射）是一个同构（换言之，\(K_{s}\) 到 E 上的每个线性映射都是同构）。

由于映射 \(\xi \mapsto \xi a\)（\(K_{s}\) 到 E 上的映射）是双射且连续（I, 第 1 页, def. 1），只需证明它是双连续的。设 \(\alpha\) 是一个实数 \(>0\)，我们证明存在 E 中 0 的邻域 V，使得若 \(\xi a \in V\)，则 \(|\xi| < \alpha\)。由于 K 不是离散的，存在元素 \(\xi_0 \in \mathbf{K}\) 使得 \(0 < |\xi_0| < \alpha\)；又由于 E 是 Hausdorff 的，存在 0 的一个邻域 V 使得 \(\xi_0a\) 不属于 V。可以假设 V 是均衡的（I, 第 7 页, prop. 4）。但这时若 \(\xi a \in \mathbf{V}\) 且 \(|\xi| \geqslant |\xi_0|\)，我们就有 \(|\xi_0\xi^{-1}| \leqslant 1\)，以及 \(\xi_0a = (\xi_0\xi^{-1})(\xi a) \in \mathbf{V}\)；由于最后这个论断不成立，我们看出 \(\xi a \in \mathbf{V}\) 蕴含 \(|\xi| < |\xi_0| < \alpha\)。证毕。

推论 1. --- 在非离散赋值除环 K 上的 Hausdorff 拓扑向量空间 E 中，每个一维向量子空间 D 都同构于 \(K_{s}\)。

推论 2. --- 设 E 是非离散赋值除环上的一个拓扑向量空间。每个作为闭齐次超平面 H 的代数补的（一维）向量子空间 D 也是 H 的拓扑补。

在 D 中，集合 \(\{0\}\) 是闭的，因为它是 D 与闭集 H 的交；因而 D 是 Hausdorff 的。但由于 E/H 也是 Hausdorff 的，D 到 E/H 上的典范映射（它是线性的）由 prop. 2 也是一个同构，由此即得结论（GT, III, § 6.2）。

定理 1. --- 设 E 是非离散赋值除环上的一个拓扑向量空间。设 H 是 E 中由方程 \(f(x) = \alpha\) 定义的超平面，其中 f 是一个不恒为零的线性型。则 H 在 E 中是闭的，当且仅当 f 连续。

该条件显然是充分的（GT, I, § 2.1, th. 1）；我们证明它是必要的。可以假设 H 是由方程 \(f(x) = 0\) 定义的闭齐次超平面。于是商空间 E/H 是 K 上的一个一维 Hausdorff 拓扑向量空间。我们可以写出 \(f = g \circ \phi\)，其中 \(\phi\) 是 E 到 E/H 上的典范映射，而 g 是 E/H 到 \(K_{s}\) 上的一个线性映射；由 prop. 2，g 连续，因此 f 也连续。

推论. --- E 上每个不恒为零的连续线性型都是 E 到 \(K_{s}\) 上的一个严格态射。

注记. --- 存在完备非离散赋值除环上的赋范拓扑向量空间的例子，其中每个连续线性型都恒为零（I, 第 25 页, 习题 4）；因此在这样的空间中，每个超平面都处处稠密（I, 第 11 页, corollary）。

